{
\renewcommand{\arraystretch}{1.3}
\setlength{\tabcolsep}{4pt}
\small
\begin{landscape}
\pagestyle{plain}
\begin{longtable}{|>{\raggedright\arraybackslash}p{1.5cm}|>{\raggedright\arraybackslash}p{1.3cm}|>{\raggedright\arraybackslash}p{3.0cm}|>{\raggedright\arraybackslash}p{2.8cm}|>{\raggedright\arraybackslash}p{4.3cm}|>{\raggedright\arraybackslash}p{2.8cm}|>{\raggedright\arraybackslash}p{4.3cm}|}
\caption{Danh sách test case E2E --- English Prep}
\label{tab:e2e-englishprep} \\
\hline
\textbf{Phân loại} & \textbf{Mã TC} & \textbf{Mô tả} & \textbf{Điều kiện tiên quyết} & \textbf{Các bước kiểm thử} & \textbf{Dữ liệu kiểm thử} & \textbf{Kết quả mong đợi} \\
\hline
\endfirsthead

\hline
\textbf{Phân loại} & \textbf{Mã TC} & \textbf{Mô tả} & \textbf{Điều kiện tiên quyết} & \textbf{Các bước kiểm thử} & \textbf{Dữ liệu kiểm thử} & \textbf{Kết quả mong đợi} \\
\hline
\endhead

\hline
\endfoot

Đăng nhập & AUTH 001 & Đăng nhập bằng email và mật khẩu thành công & Tài khoản đã được đăng ký & \begin{enumerate}[leftmargin=*, nosep, labelsep=0.3em]
    \item Mở trang đăng nhập
    \item Nhập email vào ô ''Email''
    \item Nhập mật khẩu vào ô ''Mật khẩu''
    \item Nhấn nút ''Đăng nhập''
\end{enumerate} & Email: ''test@example.com'' \newline
Mật khẩu: ''Password123!'' & Đăng nhập thành công --- chuyển đến trang chính và hiển thị tên người dùng \\
\hline
 & AUTH 002 & Đăng nhập bằng Google OAuth2 thành công & Tài khoản Google đã liên kết trong hệ thống & \begin{enumerate}[leftmargin=*, nosep, labelsep=0.3em]
    \item Mở trang đăng nhập
    \item Nhấn nút ''Đăng nhập bằng Google''
    \item Chọn tài khoản Google trên trang xác thực của Google
    \item Chờ được chuyển về ứng dụng
\end{enumerate} & Không có & Đăng nhập thành công --- chuyển đến trang chính và hiển thị tên người dùng \\
\hline
 & AUTH 003 & Phiên tự động làm mới khi access token hết hạn & Đã đăng nhập; access token sắp hết hạn & \begin{enumerate}[leftmargin=*, nosep, labelsep=0.3em]
    \item Đăng nhập thành công
    \item Đợi đến khi access token hết hạn
    \item Thực hiện thao tác bất kỳ trên ứng dụng (ví dụ: mở trang hồ sơ)
\end{enumerate} & Không có & Ứng dụng tự động làm mới phiên --- người dùng không bị gián đoạn và không cần đăng nhập lại \\
\hline
Đăng ký & AUTH 004 & Đăng ký tài khoản mới thành công &  & \begin{enumerate}[leftmargin=*, nosep, labelsep=0.3em]
    \item Mở trang đăng ký
    \item Nhập tên người dùng vào ô ''Tên người dùng''
    \item Nhập email vào ô ''Email''
    \item Nhập mật khẩu vào ô ''Mật khẩu''
    \item Nhấn nút ''Đăng ký''
    \item Mở trang hồ sơ cá nhân để kiểm tra
\end{enumerate} & Username: ''newuser'' \newline
Email: ''new@example.com'' \newline
Mật khẩu: ''StrongPass123!'' & Tài khoản được tạo thành công --- chuyển đến trang chính; trang hồ sơ hiển thị đúng username đã đăng ký \\
\hline
 & AUTH 005 & Đăng ký bằng Google OAuth2 thành công & Tài khoản Google chưa có trong hệ thống & \begin{enumerate}[leftmargin=*, nosep, labelsep=0.3em]
    \item Mở trang đăng ký
    \item Nhấn nút ''Đăng ký bằng Google''
    \item Chọn tài khoản Google trên trang xác thực
    \item Chờ được chuyển về ứng dụng
\end{enumerate} & Không có & Tài khoản mới được tạo với thông tin Google --- chuyển đến trang chính \\
\hline
Đăng xuất & AUTH 006 & Đăng xuất tất cả phiên thành công & Đã đăng nhập trên nhiều thiết bị & \begin{enumerate}[leftmargin=*, nosep, labelsep=0.3em]
    \item Mở menu tài khoản (nhấn vào avatar)
    \item Nhấn ''Đăng xuất tất cả thiết bị''
    \item Xác nhận đăng xuất
    \item Mở ứng dụng trên thiết bị/trình duyệt khác đang đăng nhập
\end{enumerate} & Không có & Tất cả phiên bị đăng xuất --- thiết bị khác yêu cầu đăng nhập lại \\
\hline
Xem hồ sơ & AUTH 007 & Xem hồ sơ cá nhân thành công & Đã đăng nhập & \begin{enumerate}[leftmargin=*, nosep, labelsep=0.3em]
    \item Nhấn vào avatar hoặc tên người dùng trên thanh điều hướng
    \item Chọn ''Hồ sơ cá nhân''
\end{enumerate} & Không có & Trang hồ sơ hiển thị đầy đủ: tên người dùng, họ tên, bio, ảnh đại diện, vai trò \\
\hline
 & AUTH 008 & Xem hồ sơ công khai người dùng khác &  & \begin{enumerate}[leftmargin=*, nosep, labelsep=0.3em]
    \item Tìm kiếm hoặc nhấn vào tên người dùng khác trong danh sách
    \item Xem trang hồ sơ công khai của người đó
\end{enumerate} & Người dùng mục tiêu: ''user123'' & Hiển thị hồ sơ giới hạn: tên người dùng, họ tên, bio, ảnh đại diện (không hiển thị thông tin nhạy cảm) \\
\hline
 & AUTH 009 & Xem hồ sơ nhiều người dùng trong danh sách & Có nhiều người dùng trong hệ thống & \begin{enumerate}[leftmargin=*, nosep, labelsep=0.3em]
    \item Mở trang kết quả tìm kiếm hoặc danh sách người dùng
    \item Kiểm tra thông tin hiển thị của từng người trong danh sách
\end{enumerate} & Không có & Mỗi người dùng hiển thị đúng: tên, ảnh đại diện, bio \\
\hline
Cập nhật hồ sơ & AUTH 010 & Cập nhật hồ sơ cá nhân thành công & Đã đăng nhập & \begin{enumerate}[leftmargin=*, nosep, labelsep=0.3em]
    \item Mở trang hồ sơ cá nhân
    \item Nhấn nút ''Chỉnh sửa hồ sơ''
    \item Thay đổi tên người dùng, họ tên và bio
    \item Nhấn nút ''Lưu''
    \item Tải lại trang hồ sơ để kiểm tra
\end{enumerate} & Username: ''newname'' \newline
Họ tên: ''Nguyễn Văn A'' \newline
Bio: ''Hello world'' & Hồ sơ được cập nhật --- trang hiển thị đúng các thông tin vừa thay đổi \\
\hline
Liên kết tài khoản & AUTH 011 & Liên kết email vào tài khoản thành công & Đã đăng nhập; chưa có email liên kết & \begin{enumerate}[leftmargin=*, nosep, labelsep=0.3em]
    \item Mở trang Cài đặt tài khoản
    \item Trong mục ''Phương thức đăng nhập'', nhấn ''Liên kết email''
    \item Nhập email và mật khẩu vào form
    \item Nhấn ''Liên kết''
    \item Kiểm tra danh sách phương thức đăng nhập
\end{enumerate} & Email: ''link@example.com'' \newline
Mật khẩu: ''LinkPass123!'' & Email mới xuất hiện trong danh sách phương thức đăng nhập \\
\hline
 & AUTH 012 & Liên kết Google vào tài khoản thành công & Đã đăng nhập; chưa có Google liên kết & \begin{enumerate}[leftmargin=*, nosep, labelsep=0.3em]
    \item Mở trang Cài đặt tài khoản
    \item Trong mục ''Phương thức đăng nhập'', nhấn ''Liên kết Google''
    \item Chọn tài khoản Google trên trang xác thực
    \item Chờ được chuyển về trang Cài đặt
    \item Kiểm tra danh sách phương thức đăng nhập
\end{enumerate} & Không có & Google mới xuất hiện trong danh sách phương thức đăng nhập \\
\hline
Đổi mật khẩu & AUTH 013 & Đổi mật khẩu thành công & Đã đăng nhập; có email liên kết & \begin{enumerate}[leftmargin=*, nosep, labelsep=0.3em]
    \item Mở trang Cài đặt tài khoản
    \item Nhấn ''Đổi mật khẩu''
    \item Nhập mật khẩu mới vào form
    \item Nhấn ''Lưu''
    \item Đăng xuất
    \item Đăng nhập lại bằng mật khẩu mới
\end{enumerate} & Mật khẩu mới: ''NewStrongPass456!'' & Đổi mật khẩu thành công --- đăng nhập lại bằng mật khẩu mới thành công \\
\hline
Xóa phương thức đăng nhập & AUTH 014 & Xóa phương thức đăng nhập thành công & Đã đăng nhập; có ít nhất 2 phương thức (email + Google) & \begin{enumerate}[leftmargin=*, nosep, labelsep=0.3em]
    \item Mở trang Cài đặt tài khoản
    \item Trong mục ''Phương thức đăng nhập'', tìm phương thức muốn xóa
    \item Nhấn nút ''Xóa'' bên cạnh phương thức đó
    \item Xác nhận xóa trong hộp thoại
    \item Kiểm tra danh sách phương thức đăng nhập
\end{enumerate} & Không có & Phương thức đã bị xóa --- không còn hiển thị trong danh sách \\
\hline
 & AUTH 015 & Xem danh sách phương thức đăng nhập & Đã đăng nhập & \begin{enumerate}[leftmargin=*, nosep, labelsep=0.3em]
    \item Mở trang Cài đặt tài khoản
    \item Xem mục ''Phương thức đăng nhập''
\end{enumerate} & Không có & Hiển thị danh sách phương thức đã liên kết (Email, Google) kèm loại và trạng thái \\
\hline
Tìm kiếm người dùng & AUTH 016 & Tìm kiếm người dùng theo tên &  & \begin{enumerate}[leftmargin=*, nosep, labelsep=0.3em]
    \item Nhấn vào ô tìm kiếm trên thanh điều hướng
    \item Nhập tên người dùng cần tìm
    \item Nhấn Enter hoặc nhấn nút tìm kiếm
    \item Xem danh sách kết quả
\end{enumerate} & Từ khóa: ''test'' & Danh sách người dùng phù hợp hiển thị với phân trang --- nhấn vào để xem hồ sơ \\
\hline
 & AUTH 017 & Gợi ý người dùng khi nhập từ khóa &  & \begin{enumerate}[leftmargin=*, nosep, labelsep=0.3em]
    \item Nhấn vào ô tìm kiếm người dùng
    \item Nhập từ khóa
    \item Xem kết quả gợi ý tự động hiển thị
\end{enumerate} & Từ khóa: ''test'' & Kết quả gợi ý hiển thị nhanh khi nhập --- danh sách cập nhật theo từ khóa \\
\hline
Duyệt bài thi & EXAM 001 & Duyệt danh sách bài thi công khai &  & \begin{enumerate}[leftmargin=*, nosep, labelsep=0.3em]
    \item Mở trang ''Luyện tập'' từ thanh điều hướng
    \item Xem danh sách bài thi hiển thị
\end{enumerate} & Không có & Danh sách bài thi hiển thị với tên, thẻ phân loại, thống kê lượt làm --- có phân trang \\
\hline
 & EXAM 002 & Lọc bài thi theo tên &  & \begin{enumerate}[leftmargin=*, nosep, labelsep=0.3em]
    \item Mở trang ''Luyện tập''
    \item Nhập từ khóa vào ô tìm kiếm bài thi
    \item Xem danh sách kết quả cập nhật
\end{enumerate} & Từ khóa: ''TOEIC'' & Chỉ hiển thị bài thi có tên chứa từ khóa đã nhập \\
\hline
 & EXAM 003 & Lọc bài thi theo thẻ phân loại &  & \begin{enumerate}[leftmargin=*, nosep, labelsep=0.3em]
    \item Mở trang ''Luyện tập''
    \item Mở bộ lọc thẻ phân loại
    \item Chọn thẻ muốn lọc
    \item Xem danh sách kết quả cập nhật
\end{enumerate} & Thẻ: ''Grammar'' & Chỉ hiển thị bài thi thuộc thẻ phân loại đã chọn \\
\hline
 & EXAM 004 & Sắp xếp danh sách bài thi &  & \begin{enumerate}[leftmargin=*, nosep, labelsep=0.3em]
    \item Mở trang ''Luyện tập''
    \item Nhấn vào tùy chọn sắp xếp
    \item Chọn tiêu chí sắp xếp (ví dụ: tên A-Z)
    \item Xem danh sách cập nhật thứ tự
\end{enumerate} & Sắp xếp: tên tăng dần & Danh sách bài thi được sắp xếp lại theo tiêu chí đã chọn \\
\hline
Xem chi tiết & EXAM 005 & Xem chi tiết bài thi & Bài thi đã được xuất bản & \begin{enumerate}[leftmargin=*, nosep, labelsep=0.3em]
    \item Mở trang ''Luyện tập''
    \item Nhấn vào bài thi muốn xem
    \item Xem trang chi tiết bài thi
\end{enumerate} & Không có & Trang chi tiết hiển thị: cấu trúc các phần, số câu hỏi mỗi phần, thẻ phân loại, tổng thời gian, tổng lượt làm \\
\hline
 & EXAM 006 & Xem thống kê bài thi & Bài thi đã được xuất bản & \begin{enumerate}[leftmargin=*, nosep, labelsep=0.3em]
    \item Mở trang chi tiết bài thi
    \item Chọn tab ''Thống kê''
\end{enumerate} & Không có & Hiển thị điểm trung bình, thời gian làm bài trung bình, tỷ lệ đúng theo từng thẻ phân loại \\
\hline
 & EXAM 007 & Xem chi tiết câu hỏi trong bài xem lại & Đã nộp bài và đang xem lại kết quả & \begin{enumerate}[leftmargin=*, nosep, labelsep=0.3em]
    \item Mở trang xem lại kết quả bài thi
    \item Nhấn vào câu hỏi cần xem chi tiết
\end{enumerate} & Không có & Hiển thị nội dung câu hỏi, các lựa chọn, giải thích đáp án, tệp đính kèm (nếu có) \\
\hline
Làm bài thi & EXAM 008 & Bắt đầu làm bài thi thành công & Đã đăng nhập; bài thi đã xuất bản & \begin{enumerate}[leftmargin=*, nosep, labelsep=0.3em]
    \item Mở trang chi tiết bài thi
    \item Nhấn nút ''Bắt đầu làm bài''
    \item Xác nhận bắt đầu trong hộp thoại
\end{enumerate} & Không có & Chuyển đến giao diện làm bài --- hiển thị câu hỏi đầu tiên và bộ đếm thời gian \\
\hline
 & EXAM 009 & Bắt đầu làm bài với tùy chọn phần và thời gian & Đã đăng nhập; bài thi có nhiều phần & \begin{enumerate}[leftmargin=*, nosep, labelsep=0.3em]
    \item Mở trang chi tiết bài thi
    \item Nhấn nút ''Bắt đầu làm bài''
    \item Bỏ chọn các phần không muốn làm
    \item Nhập thời gian tùy chỉnh
    \item Nhấn ''Bắt đầu''
\end{enumerate} & Phần chọn: ''Reading'' \newline
Thời gian: 30 phút & Giao diện làm bài chỉ chứa câu hỏi từ phần đã chọn, bộ đếm thời gian hiển thị đúng thời gian tùy chỉnh \\
\hline
 & EXAM 010 & Chọn đáp án cho câu hỏi & Đã bắt đầu lượt làm bài & \begin{enumerate}[leftmargin=*, nosep, labelsep=0.3em]
    \item Đọc nội dung câu hỏi trong giao diện làm bài
    \item Nhấn chọn đáp án (A/B/C/D)
    \item Chuyển sang câu tiếp theo bằng nút ''Câu sau''
    \item Quay lại câu trước để kiểm tra đáp án đã lưu
\end{enumerate} & Đáp án chọn: A & Đáp án được lưu --- khi quay lại câu hỏi vẫn hiển thị đáp án đã chọn \\
\hline
 & EXAM 011 & Xóa đáp án đã chọn & Đã chọn đáp án cho câu hỏi & \begin{enumerate}[leftmargin=*, nosep, labelsep=0.3em]
    \item Quay lại câu hỏi đã trả lời
    \item Nhấn ''Xóa đáp án'' hoặc bỏ chọn đáp án hiện tại
\end{enumerate} & Không có & Đáp án bị xóa --- câu hỏi trở về trạng thái chưa trả lời trên bảng câu hỏi \\
\hline
 & EXAM 012 & Đánh dấu câu hỏi để xem lại & Đã bắt đầu lượt làm bài & \begin{enumerate}[leftmargin=*, nosep, labelsep=0.3em]
    \item Trong giao diện làm bài, nhấn biểu tượng ''Đánh dấu'' (cờ/ngôi sao) trên câu hỏi
    \item Mở bảng tổng quan câu hỏi để kiểm tra
\end{enumerate} & Không có & Câu hỏi được đánh dấu --- hiển thị biểu tượng đặc biệt trên bảng tổng quan câu hỏi \\
\hline
 & EXAM 013 & Thêm ghi chú cho câu hỏi & Đã bắt đầu lượt làm bài & \begin{enumerate}[leftmargin=*, nosep, labelsep=0.3em]
    \item Trong giao diện làm bài, nhấn nút ''Ghi chú'' trên câu hỏi
    \item Nhập nội dung ghi chú vào ô văn bản
    \item Nhấn ''Lưu ghi chú''
\end{enumerate} & Ghi chú: ''Cần xem lại ngữ pháp'' & Ghi chú được lưu --- hiển thị khi mở lại câu hỏi \\
\hline
 & EXAM 014 & Xem bảng tổng quan trạng thái bài làm & Đã trả lời vài câu và đánh dấu vài câu & \begin{enumerate}[leftmargin=*, nosep, labelsep=0.3em]
    \item Trong giao diện làm bài, nhấn mở bảng tổng quan câu hỏi
    \item Kiểm tra trạng thái hiển thị của từng câu
\end{enumerate} & Không có & Bảng hiển thị đúng trạng thái: câu đã trả lời, câu đánh dấu, câu chưa làm, thời gian còn lại \\
\hline
 & EXAM 015 & Nộp bài thi thành công & Đã bắt đầu lượt làm bài và trả lời câu hỏi & \begin{enumerate}[leftmargin=*, nosep, labelsep=0.3em]
    \item Nhấn nút ''Nộp bài''
    \item Xác nhận nộp trong hộp thoại cảnh báo
    \item Chờ hệ thống chấm điểm
\end{enumerate} & Không có & Bài thi được nộp --- chuyển đến trang kết quả hiển thị điểm số và tổng kết \\
\hline
Xem kết quả & EXAM 016 & Xem kết quả lượt làm bài đã nộp & Đã nộp bài thành công & \begin{enumerate}[leftmargin=*, nosep, labelsep=0.3em]
    \item Sau khi nộp bài, xem trang kết quả tự động hiển thị
    \item Cuộn qua từng câu hỏi để xem chi tiết kết quả
\end{enumerate} & Không có & Hiển thị kết quả từng câu: đúng/sai, đáp án đã chọn vs đáp án đúng, bảng điểm tổng kết theo phần \\
\hline
Lịch sử \& Thống kê & EXAM 017 & Xem lịch sử làm bài & Đã đăng nhập; đã hoàn thành ít nhất 1 bài thi & \begin{enumerate}[leftmargin=*, nosep, labelsep=0.3em]
    \item Mở menu người dùng
    \item Chọn ''Lịch sử làm bài''
    \item Xem danh sách các lượt làm bài
\end{enumerate} & Không có & Danh sách lượt làm bài hiển thị với: tên bài thi, điểm số, ngày làm --- có phân trang \\
\hline
 & EXAM 018 & Lọc lịch sử theo bài thi cụ thể & Đã đăng nhập; có lịch sử nhiều bài thi & \begin{enumerate}[leftmargin=*, nosep, labelsep=0.3em]
    \item Mở trang ''Lịch sử làm bài''
    \item Mở bộ lọc bài thi
    \item Chọn bài thi cụ thể muốn xem
    \item Xem danh sách cập nhật
\end{enumerate} & Bài thi: ''TOEIC Practice Test 1'' & Chỉ hiển thị các lượt làm bài của bài thi đã chọn \\
\hline
 & EXAM 019 & Xem lịch hoạt động học tập & Đã đăng nhập & \begin{enumerate}[leftmargin=*, nosep, labelsep=0.3em]
    \item Mở trang ''Thống kê cá nhân''
    \item Xem phần lịch hoạt động
    \item Nhấn vào ngày cụ thể để xem chi tiết
\end{enumerate} & Không có & Lịch hiển thị các ngày có hoạt động với số lượt làm bài --- ngày có hoạt động được tô màu \\
\hline
 & EXAM 020 & Xem thống kê cá nhân tổng hợp & Đã đăng nhập; đã làm bài thi & \begin{enumerate}[leftmargin=*, nosep, labelsep=0.3em]
    \item Mở trang ''Thống kê cá nhân''
    \item Xem các chỉ số tổng hợp
\end{enumerate} & Không có & Hiển thị: tổng số lượt làm bài, điểm trung bình, tỷ lệ đúng phân theo từng thẻ phân loại \\
\hline
Mục tiêu học tập & EXAM 021 & Tạo mục tiêu học tập thành công & Đã đăng nhập; chưa có mục tiêu & \begin{enumerate}[leftmargin=*, nosep, labelsep=0.3em]
    \item Mở trang ''Mục tiêu học tập''
    \item Nhấn ''Tạo mục tiêu mới''
    \item Chọn loại mục tiêu (số lượt làm bài / điểm số)
    \item Nhập giá trị mục tiêu và ngày hoàn thành
    \item Nhấn ''Lưu''
\end{enumerate} & Loại: số lượt làm bài \newline
Mục tiêu: 10 \newline
Ngày: 01/06/2026 & Mục tiêu được tạo --- hiển thị trên trang với thanh tiến độ \\
\hline
 & EXAM 022 & Cập nhật mục tiêu học tập & Đã đăng nhập; đã có mục tiêu & \begin{enumerate}[leftmargin=*, nosep, labelsep=0.3em]
    \item Mở trang ''Mục tiêu học tập''
    \item Nhấn ''Chỉnh sửa'' trên mục tiêu hiện tại
    \item Thay đổi giá trị mục tiêu
    \item Nhấn ''Lưu''
\end{enumerate} & Mục tiêu mới: 20 \newline
Loại: điểm số & Mục tiêu được cập nhật --- hiển thị giá trị mới và thanh tiến độ cập nhật \\
\hline
 & EXAM 023 & Xóa mục tiêu học tập & Đã đăng nhập; đã có mục tiêu & \begin{enumerate}[leftmargin=*, nosep, labelsep=0.3em]
    \item Mở trang ''Mục tiêu học tập''
    \item Nhấn ''Xóa'' trên mục tiêu
    \item Xác nhận xóa trong hộp thoại
\end{enumerate} & Không có & Mục tiêu đã bị xóa --- không còn hiển thị trên trang \\
\hline
Tạo bài thi & MGMT 001 & Tạo bài thi mới thành công & Đã đăng nhập với vai trò Mod hoặc Admin & \begin{enumerate}[leftmargin=*, nosep, labelsep=0.3em]
    \item Mở trang ''Quản lý bài thi''
    \item Nhấn nút ''Tạo bài thi mới''
    \item Nhập tiêu đề, mô tả và thời gian làm bài
    \item Nhấn ''Tạo''
    \item Kiểm tra bài thi mới trong danh sách
\end{enumerate} & Tiêu đề: ''TOEIC Practice Test 1'' \newline
Mô tả: ''Bài thi thử TOEIC'' \newline
Thời gian: 120 phút & Bài thi mới xuất hiện trong danh sách quản lý ở trạng thái ''''Nháp'''' \\
\hline
 & MGMT 002 & Xem danh sách bài thi quản lý & Đã đăng nhập Mod/Admin; có bài thi đã tạo & \begin{enumerate}[leftmargin=*, nosep, labelsep=0.3em]
    \item Mở trang ''Quản lý bài thi''
    \item Xem danh sách bài thi hiển thị
    \item Sử dụng bộ lọc theo tiêu đề, người tạo hoặc trạng thái
\end{enumerate} & Không có & Danh sách bài thi hiển thị với: tiêu đề, người tạo, trạng thái, ngày tạo --- có bộ lọc và phân trang \\
\hline
 & MGMT 003 & Xem chi tiết bài thi trong trang quản lý & Đã đăng nhập Mod/Admin; bài thi tồn tại & \begin{enumerate}[leftmargin=*, nosep, labelsep=0.3em]
    \item Mở trang ''Quản lý bài thi''
    \item Nhấn vào bài thi cần xem chi tiết
\end{enumerate} & Không có & Trang chi tiết hiển thị: trạng thái, mô tả, thời gian, danh sách các phần thi, thẻ phân loại \\
\hline
Chỉnh sửa bài thi & MGMT 004 & Chỉnh sửa thông tin bài thi thành công & Đã đăng nhập Mod/Admin; bài thi tồn tại & \begin{enumerate}[leftmargin=*, nosep, labelsep=0.3em]
    \item Mở trang chi tiết bài thi trong mục quản lý
    \item Nhấn nút ''Chỉnh sửa''
    \item Sửa tiêu đề, mô tả, thời gian
    \item Thêm thẻ phân loại mới
    \item Nhấn ''Lưu''
    \item Kiểm tra trang chi tiết hiển thị đúng
\end{enumerate} & Tiêu đề mới: ''Tên mới'' \newline
Thời gian: 60 phút \newline
Thẻ thêm: ''Grammar'' & Thông tin bài thi được cập nhật --- thẻ mới xuất hiện và các trường hiển thị giá trị mới \\
\hline
Xóa bài thi & MGMT 005 & Xóa bài thi thành công & Đã đăng nhập Mod/Admin; bài thi tồn tại & \begin{enumerate}[leftmargin=*, nosep, labelsep=0.3em]
    \item Mở trang ''Quản lý bài thi''
    \item Nhấn nút ''Xóa'' trên bài thi cần xóa
    \item Xác nhận xóa trong hộp thoại
    \item Kiểm tra danh sách bài thi
\end{enumerate} & Không có & Bài thi không còn xuất hiện trong danh sách quản lý \\
\hline
Phê duyệt bài thi & MGMT 006 & Phê duyệt và xuất bản bài thi thành công & Đã đăng nhập Admin có quyền EXAM\_APPROVE; bài thi ở trạng thái nháp & \begin{enumerate}[leftmargin=*, nosep, labelsep=0.3em]
    \item Mở trang ''Quản lý bài thi''
    \item Chọn bài thi ở trạng thái ''Nháp''
    \item Nhấn nút ''Phê duyệt'' hoặc ''Xuất bản''
    \item Mở trang ''Luyện tập'' công khai để kiểm tra
\end{enumerate} & Không có & Trạng thái bài thi chuyển thành ''Đã xuất bản'' --- bài thi hiển thị trong trang luyện tập công khai \\
\hline
Quản lý phần thi & MGMT 007 & Tạo phần thi mới trong bài thi & Đã đăng nhập Mod/Admin; bài thi tồn tại & \begin{enumerate}[leftmargin=*, nosep, labelsep=0.3em]
    \item Mở trang chi tiết bài thi trong mục quản lý
    \item Nhấn nút ''Thêm phần thi''
    \item Nhập tên phần thi
    \item Nhấn ''Tạo''
\end{enumerate} & Không có & Phần thi mới xuất hiện trong cấu trúc bài thi \\
\hline
 & MGMT 008 & Tạo phần thi con lồng bên trong phần cha & Đã đăng nhập Mod/Admin; đã có phần thi cha & \begin{enumerate}[leftmargin=*, nosep, labelsep=0.3em]
    \item Mở phần thi cha trong cấu trúc bài thi
    \item Nhấn ''Thêm phần con''
    \item Nhập tên phần con
    \item Nhấn ''Tạo''
\end{enumerate} & Không có & Phần thi con xuất hiện lồng bên trong phần cha trên giao diện \\
\hline
 & MGMT 009 & Cập nhật thông tin phần thi & Đã đăng nhập Mod/Admin; phần thi tồn tại & \begin{enumerate}[leftmargin=*, nosep, labelsep=0.3em]
    \item Nhấn nút ''Chỉnh sửa'' trên phần thi
    \item Sửa tên, chỉ dẫn, loại nội dung
    \item Thêm thẻ phân loại
    \item Nhấn ''Lưu''
\end{enumerate} & Tên: ''Reading Comprehension'' \newline
Loại nội dung: passage & Phần thi hiển thị đúng các thông tin đã cập nhật \\
\hline
 & MGMT 010 & Xóa phần thi & Đã đăng nhập Mod/Admin; phần thi tồn tại & \begin{enumerate}[leftmargin=*, nosep, labelsep=0.3em]
    \item Nhấn nút ''Xóa'' trên phần thi cần xóa
    \item Xác nhận xóa trong hộp thoại
    \item Kiểm tra cấu trúc bài thi
\end{enumerate} & Không có & Phần thi và toàn bộ câu hỏi bên trong đã bị xóa khỏi cấu trúc \\
\hline
 & MGMT 011 & Di chuyển phần thi đến vị trí mới & Đã đăng nhập Mod/Admin; có nhiều phần thi & \begin{enumerate}[leftmargin=*, nosep, labelsep=0.3em]
    \item Nhấn nút ''Di chuyển'' trên phần thi hoặc kéo thả
    \item Chọn vị trí đích và phần cha mới (nếu cần)
    \item Xác nhận di chuyển
\end{enumerate} & Không có & Phần thi hiển thị tại vị trí mới trong cấu trúc bài thi \\
\hline
Quản lý câu hỏi & MGMT 012 & Tạo câu hỏi mới trong phần thi & Đã đăng nhập Mod/Admin; phần thi tồn tại & \begin{enumerate}[leftmargin=*, nosep, labelsep=0.3em]
    \item Mở phần thi trong cấu trúc bài
    \item Nhấn nút ''Thêm câu hỏi''
\end{enumerate} & Không có & Câu hỏi trống được tạo trong phần thi --- sẵn sàng chỉnh sửa nội dung \\
\hline
 & MGMT 013 & Chỉnh sửa nội dung và đáp án câu hỏi & Đã đăng nhập Mod/Admin; câu hỏi tồn tại & \begin{enumerate}[leftmargin=*, nosep, labelsep=0.3em]
    \item Nhấn vào câu hỏi cần sửa
    \item Nhập nội dung câu hỏi và giải thích
    \item Nhập điểm số
    \item Thêm các lựa chọn đáp án và đánh dấu đáp án đúng
    \item Nhấn ''Lưu''
    \item Mở lại câu hỏi để kiểm tra
\end{enumerate} & Nội dung: ''What is...?'' \newline
Điểm: 1 \newline
Loại: trắc nghiệm một đáp án \newline
Lựa chọn: A (đúng), B, C, D & Câu hỏi hiển thị đúng nội dung và các lựa chọn --- đáp án đúng được đánh dấu \\
\hline
 & MGMT 014 & Xóa câu hỏi khỏi phần thi & Đã đăng nhập Mod/Admin; câu hỏi tồn tại & \begin{enumerate}[leftmargin=*, nosep, labelsep=0.3em]
    \item Nhấn nút ''Xóa'' trên câu hỏi cần xóa
    \item Xác nhận xóa trong hộp thoại
    \item Kiểm tra danh sách câu hỏi trong phần thi
\end{enumerate} & Không có & Câu hỏi không còn hiển thị trong phần thi \\
\hline
 & MGMT 015 & Di chuyển câu hỏi sang phần khác & Đã đăng nhập Mod/Admin; có nhiều phần và câu hỏi & \begin{enumerate}[leftmargin=*, nosep, labelsep=0.3em]
    \item Nhấn nút ''Di chuyển'' trên câu hỏi
    \item Chọn phần thi đích từ danh sách
    \item Xác nhận di chuyển
    \item Mở phần thi đích để kiểm tra
\end{enumerate} & Không có & Câu hỏi xuất hiện trong phần thi mới --- không còn trong phần cũ \\
\hline
Quản lý thẻ phân loại & MGMT 016 & Xem cây thẻ phân loại &  & \begin{enumerate}[leftmargin=*, nosep, labelsep=0.3em]
    \item Mở trang ''Quản lý thẻ phân loại''
    \item Xem cây thẻ hiển thị dạng phân cấp
\end{enumerate} & Không có & Cây thẻ hiển thị đầy đủ với các cấp lồng nhau: thẻ cha, thẻ con \\
\hline
 & MGMT 017 & Xem danh sách thẻ phân loại dạng phẳng &  & \begin{enumerate}[leftmargin=*, nosep, labelsep=0.3em]
    \item Mở trang ''Quản lý thẻ phân loại''
    \item Chuyển sang chế độ xem danh sách phẳng
\end{enumerate} & Không có & Hiển thị tất cả thẻ dạng danh sách phẳng kèm tên thẻ cha tương ứng \\
\hline
 & MGMT 018 & Tạo thẻ phân loại mới & Đã đăng nhập Mod/Admin & \begin{enumerate}[leftmargin=*, nosep, labelsep=0.3em]
    \item Mở trang ''Quản lý thẻ phân loại''
    \item Nhấn ''Thêm thẻ mới''
    \item Nhập tên thẻ
    \item Chọn thẻ cha từ cây (nếu cần)
    \item Nhấn ''Tạo''
\end{enumerate} & Tên thẻ: ''Grammar'' \newline
Thẻ cha: ''English Skills'' & Thẻ mới xuất hiện trong cây dưới thẻ cha đã chọn \\
\hline
 & MGMT 019 & Đổi tên thẻ phân loại & Đã đăng nhập Mod/Admin; thẻ tồn tại & \begin{enumerate}[leftmargin=*, nosep, labelsep=0.3em]
    \item Nhấn nút ''Chỉnh sửa'' trên thẻ cần đổi tên
    \item Nhập tên mới
    \item Nhấn ''Lưu''
\end{enumerate} & Tên mới: ''Advanced Grammar'' & Tên thẻ được cập nhật trong cây phân loại \\
\hline
 & MGMT 020 & Di chuyển thẻ phân loại đến vị trí cha mới & Đã đăng nhập Mod/Admin; thẻ tồn tại & \begin{enumerate}[leftmargin=*, nosep, labelsep=0.3em]
    \item Nhấn nút ''Di chuyển'' trên thẻ cần di chuyển
    \item Chọn thẻ cha mới từ cây
    \item Xác nhận di chuyển
\end{enumerate} & Không có & Thẻ hiển thị tại vị trí cha mới trong cây phân loại \\
\hline
 & MGMT 021 & Xóa thẻ phân loại & Đã đăng nhập Mod/Admin; thẻ tồn tại & \begin{enumerate}[leftmargin=*, nosep, labelsep=0.3em]
    \item Nhấn nút ''Xóa'' trên thẻ cần xóa
    \item Xác nhận xóa trong hộp thoại
    \item Kiểm tra cây thẻ phân loại
\end{enumerate} & Không có & Thẻ không còn xuất hiện trong cây phân loại \\
\hline
Tải ảnh công khai & FILE 001 & Tải ảnh đại diện lên thành công & Đã đăng nhập & \begin{enumerate}[leftmargin=*, nosep, labelsep=0.3em]
    \item Mở trang chỉnh sửa hồ sơ
    \item Nhấn vào vùng ảnh đại diện hoặc nút ''Chọn ảnh''
    \item Chọn tệp ảnh từ máy tính
    \item Chờ ảnh tải lên hoàn tất
    \item Nhấn ''Lưu'' hồ sơ
\end{enumerate} & Tệp: avatar.png \newline
Kích thước: 100KB \newline
Loại: image/png & Ảnh được tải lên thành công --- hiển thị xem trước trên trang hồ sơ \\
\hline
 & FILE 002 & Đính kèm tệp riêng tư vào nội dung & Đã đăng nhập & \begin{enumerate}[leftmargin=*, nosep, labelsep=0.3em]
    \item Mở form tạo/chỉnh sửa nội dung (ví dụ: câu hỏi bài thi)
    \item Nhấn nút ''Đính kèm tệp''
    \item Chọn tệp từ máy tính
    \item Chờ tệp tải lên hoàn tất
\end{enumerate} & Tệp: document.pdf \newline
Kích thước: 200KB & Tệp được đính kèm thành công --- chỉ người có quyền mới truy cập được \\
\hline
 & FILE 003 & Ảnh công khai hiển thị đúng sau khi tải lên & Đã tải ảnh lên thành công ở FILE 001 & \begin{enumerate}[leftmargin=*, nosep, labelsep=0.3em]
    \item Mở trang hồ sơ cá nhân
    \item Kiểm tra ảnh đại diện hiển thị
    \item Nhấn chuột phải vào ảnh, chọn ''Mở ảnh trong tab mới''
\end{enumerate} & Không có & Ảnh hiển thị đúng trên trang hồ sơ và có thể mở trực tiếp trong tab mới \\
\hline
Throttle Guard & RATE 001 & Bị chặn khi thao tác quá nhanh liên tục &  & \begin{enumerate}[leftmargin=*, nosep, labelsep=0.3em]
    \item Thực hiện thao tác liên tục rất nhanh (ví dụ: tải lại trang, nhấn nút liên tục hơn 60 lần trong 60 giây)
    \item Quan sát phản hồi sau lần thao tác thứ 61
\end{enumerate} & Không có & Hiển thị thông báo lỗi ''Quá nhiều yêu cầu, vui lòng thử lại sau'' --- thao tác bị từ chối \\
\hline
 & RATE 002 & Thao tác trở lại bình thường sau thời gian chờ & Đã bị giới hạn tốc độ ở RATE 001 & \begin{enumerate}[leftmargin=*, nosep, labelsep=0.3em]
    \item Đợi 60 giây sau khi bị giới hạn
    \item Thực hiện lại thao tác bình thường (ví dụ: mở trang luyện tập)
\end{enumerate} & Không có & Trang hoạt động bình thường --- nội dung hiển thị đúng không còn lỗi \\
\hline
JWT Guard & RATE 003 & Truy cập trang yêu cầu đăng nhập khi chưa đăng nhập &  & \begin{enumerate}[leftmargin=*, nosep, labelsep=0.3em]
    \item Mở trình duyệt mới (chưa đăng nhập)
    \item Truy cập trực tiếp trang hồ sơ cá nhân qua URL
\end{enumerate} & Không có & Chuyển hướng về trang đăng nhập hoặc hiển thị thông báo yêu cầu đăng nhập \\
\hline
 & RATE 004 & Phiên hết hạn khi không hoạt động lâu & Đã đăng nhập & \begin{enumerate}[leftmargin=*, nosep, labelsep=0.3em]
    \item Đăng nhập thành công
    \item Không thao tác trong thời gian dài để phiên hết hạn
    \item Thực hiện thao tác yêu cầu xác thực (ví dụ: mở hồ sơ cá nhân)
\end{enumerate} & Không có & Hệ thống tự động làm mới phiên hoặc yêu cầu đăng nhập lại \\
\hline
 & RATE 005 & Phiên bị từ chối khi token bị giả mạo &  & \begin{enumerate}[leftmargin=*, nosep, labelsep=0.3em]
    \item Sử dụng công cụ phát triển trình duyệt để sửa đổi token thành chuỗi không hợp lệ
    \item Thực hiện thao tác yêu cầu xác thực (ví dụ: mở hồ sơ cá nhân)
\end{enumerate} & Không có & Chuyển hướng về trang đăng nhập --- phiên không hợp lệ bị từ chối \\
\hline
\end{longtable}
\end{landscape}
}
